% !PRINTED 773
\weberpage{773}
%<<CONT
so sind die $c_1, c_2, \ldots$ \gesperrt{ganze Zahlen}, weil $\chi(x)$ ein Theiler von $x^m - 1$ sein muss (nach dem \gesperrt{Gauss}'schen Satze, Bd.~I, \S.\,2).

Daraus ergiebt sich, dass, wenn $\varphi(x)$ eine ganze Function von $x$ mit ganzzahligen rationalen Coëfficienten bedeutet, auch wenn der Grad von $\varphi(x)$ ursprünglich höher als $\nu$ ist, die Coëfficienten $m_0, m_1, \ldots, m_{\nu-1}$ in (9) ganze Zahlen werden; denn diese Coëfficienten erhält man, wenn man den Rest der Division von $\varphi(x)$ durch $\chi(x)$ aufsucht.

Wir nehmen nun an, $f_n(x)$ sei im Körper $R(\alpha)$ zerlegbar, und es sei demnach:

\weq{11}{f_n(x) = \varphi(x,\alpha)\ \psi(x,\alpha),}

worin $\varphi(x,\alpha)$ und $\psi(x,\alpha)$ zwei ganze Functionen von $x$ sind, deren keine constant ist, deren Coëfficienten in $R(\alpha)$ enthalten sind. Aus (8) und (11) ergiebt sich, wenn wir $x=1$ setzen,

\weq{12}{p = \varphi(1,\alpha)\ \psi(1,\alpha),}

worin $\varphi(1,\alpha)$, $\psi(1,\alpha)$ als Zahlen des Körpers $R(\alpha)$ in der Form (9) darstellbar sind. Bezeichnen wir die kleinsten gemeinschaftlichen Nenner dieser beiden Darstellungen mit $a$ und $b$, so erhalten wir:

\weq{13}{\begin{aligned}
a\,\varphi(1,\alpha) &= a_0 + a_1\alpha + a_2\alpha^2 + \cdots + a_{\nu-1}\alpha^{\nu-1} = A(\alpha)\\
b\,\psi(1,\alpha) &= b_0 + b_1\alpha + b_2\alpha^2 + \cdots + b_{\nu-1}\alpha^{\nu-1} = B(\alpha),
\end{aligned}}

worin $a, a_0, a_1, \ldots, a_{\nu-1}$ ganze Zahlen ohne gemeinsamen Theiler und $A$ ein Functionszeichen ist. Gleiches gilt für $b, b_0, b_1, \ldots, b_{\nu-1}$ und $B$.

Die Wurzeln der Gleichung $\varphi(x,\alpha) = 0$ finden sich unter den primitiven $n^{\text{ten}}$ Einheitswurzeln, und diese sind sämmtlich Potenzen von einer unter ihnen, $r$. Demnach giebt es gewisse Potenzen $r^h, r^{h'}, r^{h''}, \ldots$, so dass

\[
\varphi(x,\alpha) = (x - r^h)(x - r^{h'})(x - r^{h''}) \ldots
\]

wird, und folglich nach (13):

\[
A(\alpha) = a(1 - r^h)(1 - r^{h'})(1 - r^{h''}) \ldots
\]

Nehmen wir hiervon die $n^{\text{te}}$ Potenz, so ergiebt sich nach (1), mit Rücksicht auf die Gleichung $r^n = 1$:

\[
[A(\alpha)]^n = p\,\chi(r),
\]

worin $\chi(x)$ eine ganze Function der Variablen $x$ mit ganzzahligen Coëfficienten bedeutet.
